\paragraph{JF}
Järnvägsföretag.
De företag som innehar ett av Transportstyrelsen utfärdat trafikeringstillstånd på järnvägsinfrastruktur hos en infrastrukturförvaltare (t ex Trafikverket)~\cite{Jarnvagsforetag2023}.
\paragraph{Körplan}
Villkorat dokument för att en tågfärd ska får lov att genomföras.
Innehåller en unik identifierande beteckning som tilldelas av vederbörande infrastrukturförvaltare.
Anger bland annat vid vilken driftplats\footnote{Ett från linjen avgränsat område av banan som kan övervakas av tågklarerare mer detaljerat än vad som krävs för linjen~\cite{TTJ2015}.} eller driftplatsdel\footnote{Avgränsad del av en driftplats ~\cite{TTJ2015}.} tågfärden påbörjas respektive avslutas, samt vid vilka trafikplatser\footnote{Gemensam term för driftplats, driftplatsdel, linjeplats, hållplats och hållställe\cite{TTJ2015}.} tågfärden passerar och vid vilka tidpunkter.
~\cite{TTJ2015}
\paragraph{Tågklarerare}
Förkortas TKL. Person som från en tågtrafikledningscentral övervakar och leder trafikverksamheter på järnväg.
\paragraph{SFERA}
Smart communications for efficient rail activities.
Ett protokoll för kommunikation mellan TMS och C-DAS.
För närvarande under införande i Sverige.
\cite{Joborn2023}
\paragraph{Målpunkt}
Definitionen av en målpunkt enligt SFERA.
\paragraph{Digital graf}
Här lyder en förklaring av systemet Digital graf.
\paragraph{RTTP}
Real time traffic plan.
~\cite{Joborn2023}
\paragraph{TMS}
Traffic management system.
System som används av tågklarerare för operativ planering av kapacitet på järnväg och skickar RTTP till en C-DAS-klient för verkställande när ändringar görs.
Digital graf är ett exempel på ett sådant.
~\cite{Joborn2023}
\paragraph{DAS}
Driver advisory system.
\paragraph{C-DAS}
Connected driver advisory system.
\paragraph{S-DAS}
Standalone driver advisory system.
\paragraph{GSM-R}
Global system for mobile communications - railway.
Tillägg till den internationella mobiltelefonistandarden GSM och som specifikt avser järnvägskommunikation.
GSM-R har ett reserverat frekvensspektrum som bara får användas inom järnvägsverksamhet.
Via antenner längs med järnvägen säkerställs att tåg alltid kan nås med radiokommunikation på frekvenser som är garanterat tillgängliga.
Erbjuder säker överföring av både röst (samtal) och data (passagerarinformationssystem, lastspårning, videoövervakning m.m).
~\cite{Malux2019}